%% 课程报告 / 课堂展示（清新双色学术风）
%% 编译: xelatex
\documentclass[UTF8,aspectratio=169,10pt]{ctexbeamer}
\usepackage{ctex}
\usepackage{graphicx}
\usepackage{amsmath}
\usepackage{listings}
\usepackage{tikz}
\usetikzlibrary{positioning,arrows.meta}

\setCJKmainfont{SimSun}
\setCJKsansfont{Microsoft YaHei}
\setCJKmonofont{FangSong}

\definecolor{primary}{HTML}{0F766E}
\definecolor{primaryDark}{HTML}{115E59}
\definecolor{accent}{HTML}{F59E0B}
\definecolor{ink}{HTML}{0F172A}
\definecolor{muted}{HTML}{64748B}

\mode<presentation>
\setbeamertemplate{navigation symbols}{}
\setbeamertemplate{footline}[frame number]
\setbeamercolor{structure}{fg=primary}
\setbeamercolor{frametitle}{bg=primary,fg=white}
\setbeamercolor{title}{fg=primaryDark}
\setbeamercolor{block title}{bg=primary,fg=white}
\setbeamercolor{block body}{bg=primary!8,fg=ink}
\setbeamercolor{block title alerted}{bg=accent,fg=white}
\setbeamercolor{block body alerted}{bg=accent!12,fg=ink}
\setbeamerfont{frametitle}{size=\large,series=\bfseries}
\setbeamertemplate{itemize items}{\textcolor{primary}{$\bullet$}}
\setbeamertemplate{enumerate items}{\textcolor{primary}{\insertenumlabel.}}

\title{课程名称报告}
\subtitle{第 X 次作业 / 课堂展示}
\author{姓名 · 学号}
\institute{学校 · 学院}
\date{\today}

\begin{document}

\begin{frame}
  \titlepage
\end{frame}

\begin{frame}{大纲}
  \tableofcontents
\end{frame}

\section{问题引入}

\begin{frame}{为什么要学这个主题？}
  \begin{block}{核心问题}
    在给定约束下，如何获得可解释、可复现的解决方案？
  \end{block}
  \begin{itemize}
    \item 现实场景中的建模需求
    \item 课程知识与后续章节的衔接
    \item 本次报告的目标与范围
  \end{itemize}
\end{frame}

\section{基本原理}

\begin{frame}{关键概念}
  \begin{itemize}
    \item \textbf{定义}：用一句话说明对象与条件
    \item \textbf{性质}：列出 2--3 条关键性质
    \item \textbf{算法}：给出流程或伪代码框架
  \end{itemize}
  \begin{alertblock}{注意}
    符号约定与前置知识请在此处交代清楚。
  \end{alertblock}
\end{frame}

\begin{frame}{流程示意}
  \centering
  \begin{tikzpicture}[
      node distance=1.2cm and 1.6cm,
      box/.style={draw=primary, fill=primary!10, rounded corners=4pt, minimum height=1cm, minimum width=2.2cm, align=center},
      arr/.style={-{Stealth[length=3mm]}, thick, primary}
    ]
    \node[box] (a) {输入数据};
    \node[box, right=of a] (b) {预处理};
    \node[box, right=of b] (c) {模型};
    \node[box, right=of c] (d) {输出};
    \draw[arr] (a) -- (b);
    \draw[arr] (b) -- (c);
    \draw[arr] (c) -- (d);
  \end{tikzpicture}
\end{frame}

\section{示例与讨论}

\begin{frame}{简单算例}
  \begin{equation}
    \min_{x} \; c^{\top}x
    \quad \text{s.t.} \quad
    Ax \le b,\; x \ge 0
  \end{equation}
  \begin{itemize}
    \item 说明变量含义与约束解释
    \item 给出求解思路或软件命令
  \end{itemize}
\end{frame}

\begin{frame}{小结}
  \begin{enumerate}
    \item 回顾本次报告的 3 个要点
    \item 指出易错点与作业提示
    \item 预告下一节内容
  \end{enumerate}
\end{frame}

\begin{frame}{谢谢}
  \centering
  {\Large 感谢聆听，欢迎提问！}
\end{frame}

\end{document}
